\section*{Bab \ref{ch:g1:plants-around-us} --- Tumbuhan di Sekitar Kita}

\begin{solution}{exo:g1:plants-around-us:1}
\omterm{def:g1:plants-around-us:parts}{Akar} (di dalam tanah), \omterm{def:g1:plants-around-us:parts}{batang}, \omterm{def:g1:plants-around-us:parts}{daun}, \omterm{def:g1:plants-around-us:parts}{bunga}.
\end{solution}

\begin{solution}{exo:g1:plants-around-us:2}
(a) \omterm{def:g1:plants-around-us:parts}{Akarnya}. (b) \omterm{def:g1:plants-around-us:parts}{Daunnya}. (c) \omterm{def:g1:plants-around-us:parts}{Bunganya}. (d) \omterm{def:g1:plants-around-us:parts}{Batangnya}.
\end{solution}

\begin{solution}{exo:g1:plants-around-us:3}
Ia lahir dari biji, ia makan (air dari tanah, cahaya dari langit), ia
tumbuh, ia membuat biji yang menjadi \omterm{def:g1:plants-around-us:parts}{bunga} aster baru, dan ia mati
pada suatu hari. Semua tanda kehidupan ada di situ; berjalan bukan
salah satunya.
\end{solution}

\begin{solution}{exo:g1:plants-around-us:4}
\omterm{def:g1:plants-around-us:parts}{Batang} sebuah \omterm{def:g1:plants-around-us:tree}{pohon} adalah \omterm{def:g1:plants-around-us:tree}{batang pokoknya}; \omterm{def:g1:plants-around-us:parts}{daunnya} dibawa oleh
\omterm{def:g1:plants-around-us:tree}{ranting-rantingnya}.
\end{solution}

\begin{solution}{exo:g1:plants-around-us:5}
Air, cahaya, dan tempat bagi \omterm{def:g1:plants-around-us:parts}{akarnya} --- biasanya tanah.
\end{solution}

\begin{solution}{exo:g1:plants-around-us:6}
Penyiramannya. \omterm{def:g1:plants-around-us:plant}{Tumbuhannya} layu karena kekurangan air; beberapa jam
sesudah disiram, \omterm{def:g1:plants-around-us:parts}{daunnya} akan tegak kembali.
\end{solution}

\begin{solution}{exo:g1:plants-around-us:7}
Sama: keduanya \omterm{def:g1:plants-around-us:plant}{tumbuhan}, dengan \omterm{def:g1:plants-around-us:parts}{akar}, \omterm{def:g1:plants-around-us:parts}{batang}, \omterm{def:g1:plants-around-us:parts}{daun} dan \omterm{def:g1:plants-around-us:parts}{bunga}, dan
keduanya membutuhkan air dan cahaya. Berbeda: \omterm{def:g1:plants-around-us:parts}{batang} \omterm{def:g1:plants-around-us:tree}{pohonnya} adalah
\omterm{def:g1:plants-around-us:tree}{batang pokok} berkayu sedangkan \omterm{def:g1:plants-around-us:parts}{batang} \omterm{def:g1:plants-around-us:parts}{bunga} asternya lunak dan hijau,
dan \omterm{def:g1:plants-around-us:tree}{pohonnya} jauh lebih besar serta hidup jauh lebih lama.
\end{solution}

\begin{solution}{exo:g1:plants-around-us:8}
Tanaman dekat jendela seharusnya hijau dan tegak; tanaman dalam lemari
pucat, memanjang dan lemah, walaupun ia disiram. Kebutuhan yang hilang
adalah cahaya --- air saja tidak cukup.
\end{solution}

\begin{solution}{exo:g1:plants-around-us:9}
Misalnya: \omterm{def:g1:plants-around-us:plant}{tumbuhannya} tumbuh --- makin tinggi dari minggu ke minggu
--- dan ia berputar atau condong ke arah cahaya, seperti \omterm{def:g1:plants-around-us:parts}{bunga}
matahari yang mengikuti hari. Lambat tidak sama dengan tidak berbuat
apa-apa.
\end{solution}

\begin{solution}{exo:g1:plants-around-us:10}
\omterm{def:g1:plants-around-us:parts}{Akarnya}. ``Di dalam tanah'' adalah tempat \omterm{def:g1:plants-around-us:parts}{akar} pada gambarnya ---
wortel yang kita makan adalah \omterm{def:g1:plants-around-us:parts}{akar} jingga yang besar dan gemuk, yang
diisi oleh \omterm{def:g1:plants-around-us:plant}{tumbuhannya} dengan makanan.
\end{solution}
